\section{Nonexpansion and finite grids}\label{app:global}

We return to the canonical morphism $f\colon X\to Z$.  Put
$X^0=f^{-1}(Z_{\reg})$, and let $h$ be the orthogonal product metric
associated with the parallel splitting on $X^0$ from
Proposition~\ref{canon:pseudodistance}.
We use the following established properties of $h$: it is real analytic,
it is locally a product metric for $E_1\oplus E_2$, and its intrinsic
length metric has a completion identified homeomorphically with $Z$.
Write $\rho$ for the resulting distance on $Z$, and recall
\[
 d(x,x')=\rho(f(x),f(x'))\qquad(x,x'\in X).
\]
Thus $d$ is a continuous pseudodistance, whose zero-distance classes are
the fibers of $f$.  Diameters below are taken for $d$.

For each reduced fiber $P=f^{-1}(z)_{\mathrm{red}}$, Theorem~\ref{thm:crossing}
provides a connected open neighborhood $U\supset P$ and a holomorphic
crossing map $c\colon U\times U\to X$.  We will use its identities
\begin{equation}\label{global:crossing-identities}
 \begin{gathered}
 c(x,x)=x,\qquad
 \dd_1c(T_X)\subset E_1,\quad \dd_1c(E_2)=0,\\
 \dd_2c(T_X)\subset E_2,\quad \dd_2c(E_1)=0.
 \end{gathered}
\end{equation}
Near a diagonal point, $c$ is the ordinary operation that takes the first
coordinate of the first input and the second coordinate of the second.

\subsection{Nonexpansion across the exceptional locus}

\begin{lemma}\label{global:nonexpansion}
For each fiber value $z\in Z$, one can choose a crossing map as above and
$r>0$ such that $f^{-1}(B_{\rho}(z,r))\subset U$ and
\begin{equation}\label{global:lipschitz}
 \begin{split}
 d\bigl(c(a,b),c(a',b)\bigr)&\leq d(a,a'),\\
 d\bigl(c(a,b),c(a,b')\bigr)&\leq d(b,b')
 \end{split}
\end{equation}
whenever all the inputs in the relevant inequality belong to
$f^{-1}(B_{\rho}(z,r/10))$.
\end{lemma}

\begin{proof}
On the open set
\[
 V=\{(a,b)\in U\times U:a,b,c(a,b)\in X^0\}
\]
we claim the tensor identities
\begin{equation}\label{global:metric-identities}
 \begin{split}
 \bigl\lVert\dd_1c(v)\bigr\rVert_{h,c(a,b)}^2
     &=\lVert v_1\rVert_{h,a}^2,\\
 \bigl\lVert\dd_2c(w)\bigr\rVert_{h,c(a,b)}^2
     &=\lVert w_2\rVert_{h,b}^2,
 \end{split}
\end{equation}
where $v=v_1+v_2$ and $w=w_1+w_2$ denote the supplied splittings of real
tangent vectors.  The identities hold near any regular diagonal point
by the local product expression for $h$.  The complement of $V$ is a
proper complex analytic subset of the connected manifold $U\times U$,
so $V$ is connected.  Both sides of each identity are real analytic on
$V$.  The real analytic identity principle proves
\eqref{global:metric-identities} throughout $V$.

Properness of $f$ and the topology induced by $\rho$ give $r>0$ with
$f^{-1}(B_{\rho}(z,r))\subset U$.  Consider the first inequality in
\eqref{global:lipschitz}.  Approximate $(a,a',b)$ in $X^3$ by triples
$(a_m,a'_m,b_m)$ for which all three inputs and both output points are in
$X^0$.  Such triples are dense: each excluded condition is a proper
analytic subset of $U^3$, and properness of the output conditions follows
from the regular values of $c$ near a regular diagonal point.

For large $m$ all inputs lie over $B_{\rho}(z,r/9)$.  Join $a_m$ to $a'_m$
in $X^0$ by a piecewise smooth path of $h$-length less than
$d(a_m,a'_m)+\eta$, with $\eta>0$ arbitrarily small.  This is possible
because $d$ on $X^0$ is its intrinsic length distance.  Every point of
this path lies over $B_{\rho}(z,r/3+\eta)$: the initial point has distance
less than $r/9$ from $z$, and the path length is less than $2r/9+\eta$.
In particular the path remains in $U$, with a margin for perturbation.

For this fixed $b_m$, the set
\[
 A_m=\{x\in U\cap X^0:c(x,b_m)\notin X^0\}
\]
is a proper analytic subset, since the output at $a_m$ is regular.
It has real codimension at least two.  Keeping the endpoints fixed, we
may perturb the path to avoid $A_m$ with arbitrarily small length error.
Indeed, on a neighborhood of the compact path the metric is smooth;
general position relative to an analytic stratification gives such a
perturbation within that neighborhood.  Along the perturbed path,
\eqref{global:metric-identities} bounds the output speed by the input
speed, because the two summands are $h$-orthogonal.  Integration, followed
by letting the length errors tend to zero, proves
\[
 d\bigl(c(a_m,b_m),c(a'_m,b_m)\bigr)\leq d(a_m,a'_m).
\]
Continuity of $c$ and $d$ gives the first inequality for the original
inputs.  Interchanging the two variables proves the second.
\end{proof}

Compactness of $Z$ now supplies a uniform constant $\delta>0$ with the
following property:
\begin{equation}\label{global:uniform}
 \begin{gathered}
 \text{Every set }S\subset X\text{ with }\operatorname{diam}_d(S)\leq2\delta
 \text{ lies in a crossing domain }U,\\
 \text{and \eqref{global:lipschitz} holds for all inputs taken from }S.
 \end{gathered}
\end{equation}
To see this, take a finite subcover of $Z$ by the smaller balls in
Lemma~\ref{global:nonexpansion}, and choose $2\delta$ strictly below a
Lebesgue number.  No ordinary diameter bound for the fibers is asserted
or needed here.

\subsection{Two finite grid constructions}

For intervals $I,J$, the box condition of
Definition~\ref{rat:definition} says that the first product coordinate of
$R:I\times J\to X$ depends locally only on the parameter in $I$, and
the second only on the parameter in $J$. Its horizontal and vertical
paths consequently follow the $E_1$- and $E_2$-plaques, respectively.
We first record the rule for filling one tile.

Suppose two prescribed edges $\alpha\colon I\to X$ and
$\beta\colon J\to X$ meet at a corner $p$, follow the first and second
foliations respectively, and each have diameter at most $\delta$.
Their union has diameter at most $2\delta$.  Choose a crossing map by
\eqref{global:uniform} and set
\begin{equation}\label{global:tile}
 R(s,t)=c(\alpha(s),\beta(t)).
\end{equation}
This is a box and agrees with both prescribed edges.  For example, for
fixed $s$, vary the second input along the part of $\alpha$ joining $p$
to $\alpha(s)$.  The identity $\dd_2c(E_1)=0$ makes the value constant,
and therefore
\[
 c(\alpha(s),p)=c(\alpha(s),\alpha(s))=\alpha(s).
\]
The other edge follows from $\dd_1c(E_2)=0$.  This argument applies to
continuous plaque paths, since the corresponding functions are locally
constant on their parameter intervals.  Finally,
\eqref{global:lipschitz} shows that every horizontal or vertical partial
slice of this tile has diameter at most $\delta$.  Thus the same bound
is available for its two new edges.

\begin{lemma}\label{global:boxes}
The following boxes exist.
\begin{enumerate}
 \item For every continuous path $\lambda\colon[0,1]\to X$, there is a
 box $R\colon[0,1]^2\to X$ such that $R(s,s)=\lambda(s)$.
 \item For continuous paths $\alpha$ and $\beta$ in leaves of the first
 and second foliations, respectively, with
 $\alpha(0)=\beta(0)$, there is a box on $[0,1]^2$ satisfying
 $R(s,0)=\alpha(s)$ and $R(0,t)=\beta(t)$.
\end{enumerate}
Here continuity of the prescribed leaf paths is in the intrinsic leaf
topology.
\end{lemma}

\begin{proof}
For the first assertion choose a subdivision
$0=t_0<\cdots<t_N=1$ such that each $\lambda(I_i)$ has diameter at most
$\delta$, where $I_i=[t_{i-1},t_i]$.  On every diagonal tile $I_i^2$ put
\[
 R(s,t)=c_i(\lambda(s),\lambda(t)),
\]
using \eqref{global:uniform}.  Its diagonal is $\lambda$, and each of
its partial slices has diameter at most $\delta$ by
\eqref{global:lipschitz}.  These tiles agree at their common vertices.

Fill the remaining tiles $I_i\times I_j$ by increasing $|i-j|$.
Above the diagonal, the bottom and right edges of a new tile have
already been defined; below it, the top and left edges have already
been defined.  In either case these are one horizontal and one vertical
edge, meeting at a corner.  Apply \eqref{global:tile}.
The induction invariant is that every partial slice of each completed
tile has diameter at most $\delta$.  The tile rule proves both agreement
on its old edges and preservation of this invariant.  Tiles at the same
stage share no edge.  If they share a vertex, its value is already fixed
by an intervening tile from an earlier stage.  Thus the finite filling
is consistent everywhere.

For the second assertion, subdivide each prescribed axis so that the
image of every edge segment has diameter at most $\delta$.  Fill the
rectangular array in increasing order of the sum of its two indices.
The bottom and left edges of a new tile have already been prescribed,
so precisely the same rule and invariant apply.

In both constructions the tile maps glue continuously.  The box
property also holds across seams: near a seam value, continuity places
the adjacent path pieces in one foliation chart, and both follow the
plaque through their common endpoint.  This proves the assertions.
\end{proof}

The filling orders are the same as in Figure~\ref{fig:grid}. Here the
one-tile rule \eqref{global:tile} propagates the diameter bound
\(\delta\) by separate nonexpansion, rather than preserving a
holomorphic parametrization on a full polydisk.

The construction uses only continuity and diameter control.  In
particular it applies to paths inside a collapsed fiber, whose
pseudodiameter is zero, without assigning those paths any finite length
for the limiting metric.

\subsection{The resulting product of universal covers}

The two grid constructions have now supplied the exact inputs for
transport of inverse germs in Proposition~\ref{prop:germ-transport}.
That proposition uses only product charts and a continuous box; it does
not use the compact K\"ahler splitting theorem or the construction of
boxes in Sections~\ref{sec:hartogs}--\ref{sec:global}. We spell out the
remaining application to keep the alternative argument independent of
that theorem.

Fix \(o\in X\), and give the leaves \(L_1,L_2\) through \(o\)
their intrinsic complex-manifold structures, as in
Lemma~\ref{lem:leaves}. Let
\[
 H_o:(L_1\times L_2,(o,o))\longrightarrow(X,o),\qquad
 K_o:(X,o)\longrightarrow(L_1\times L_2,(o,o))
\]
be inverse local product germs. Every path in \(X\) from \(o\) is
the diagonal of a box by Lemma~\ref{global:boxes}(1). Its axes lie
continuously in the fixed intrinsic leaves, so
Proposition~\ref{prop:germ-transport} continues \(K_o\) along that
path. Conversely, a path in \(L_1\times L_2\) consists of two
intrinsically continuous leaf paths. Lemma~\ref{global:boxes}(2) supplies
a box with those axes, and the same transport proposition continues
\(H_o\) along the prescribed path. All continued germs are local
biholomorphisms.

\begin{proposition}\label{global:product}
Under the hypotheses of Corollary~\ref{thm:nef-big}, let \(Y\to X\)
be the universal cover and \(\widetilde L_i\to L_i\) the universal
covers of the leaves through \(o\). Then
\[
 Y\simeq\widetilde L_1\times\widetilde L_2
\]
biholomorphically, with the lifts of \(E_i\) identified with the
respective tangent bundles of the factors.
\end{proposition}

\begin{proof}
All-path continuation and monodromy give holomorphic maps
\[
 F:Y\longrightarrow L_1\times L_2,
 \qquad G:\widetilde L_1\times\widetilde L_2\longrightarrow X.
\]
With compatible basepoints, the covering lifting theorem gives
holomorphic lifts
\[
 \widehat F:Y\longrightarrow\widetilde L_1\times\widetilde L_2,
 \qquad \widehat G:\widetilde L_1\times\widetilde L_2\longrightarrow Y.
\]
Their initial germs are inverse. Each composite is therefore the
identity near its basepoint, and hence everywhere on its connected
source by the identity theorem. The wrong-factor components of their
differentials are holomorphic bundle morphisms that vanish in the
initial product chart. They vanish identically, giving the asserted
compatibility with the supplied splitting.
\end{proof}

Proposition~\ref{global:product} completes the alternative proof of
Corollary~\ref{thm:nef-big}. Its continuation mechanism is the separate
nonexpansion of crossings around collapsed canonical fibres together
with the finite-grid diameter bounds.
